\section{Phase-Transformation Relationship Construction From Plane-Direction Correspondence}

PyTex supports a stable constructor for an orientation relationship when the user can state one
parent-plane / child-plane correspondence together with one parent-direction / child-direction
correspondence.

\subsection{Defining Data}

The constructor consumes:

\begin{itemize}
  \item a parent-plane normal in the parent crystal frame
  \item a child-plane normal in the child crystal frame
  \item a parent direction in the parent crystal frame
  \item a child direction in the child crystal frame
\end{itemize}

The direction inputs are interpreted as in-plane guidance vectors. PyTex projects them into the
corresponding plane before constructing the final basis, so the stable contract is parallelism, not
pre-normalized basis coordinates.

\subsection{Construction Rule}

For each phase, PyTex builds a right-handed orthonormal frame:

\begin{align}
\hat{\mathbf{x}} &= \mathrm{normalize}\left(\mathbf{d} - (\mathbf{d}\cdot\hat{\mathbf{n}})\hat{\mathbf{n}}\right) \\
\hat{\mathbf{y}} &= \mathrm{normalize}\left(\hat{\mathbf{n}} \times \hat{\mathbf{x}}\right) \\
\hat{\mathbf{z}} &= \hat{\mathbf{n}}
\end{align}

where $\hat{\mathbf{n}}$ is the unit plane normal and $\mathbf{d}$ is the supplied direction.

Let $F_p$ be the parent frame assembled from the parent correspondence and $F_c$ be the child
frame assembled from the child correspondence. The parent-to-child rotation returned by PyTex is

\begin{equation}
R_{p \rightarrow c} = F_c F_p^\mathsf{T}.
\end{equation}

This is the same basis-construction doctrine already used for orientation construction from plane
and direction correspondence elsewhere in the library, but here both sides are crystal frames
rather than crystal-to-specimen mapping.

\subsection{Assumptions And Failure Modes}

\begin{itemize}
  \item The parent plane and parent direction must belong to the same parent phase.
  \item The child plane and child direction must belong to the same child phase.
  \item The projected direction must not become degenerate after projection into the plane.
  \item The constructor currently treats one plane plus one direction as the defining evidence. It
  does not yet reconcile redundant, noisy, or conflicting correspondence sets.
\end{itemize}

\subsection{Interpretation}

The resulting rotation is a crystal-to-crystal mapping between the parent and child phases. It is
not, by itself, a full parent reconstruction algorithm or a named literature catalog. Those broader
surfaces remain separate concerns.

\subsection{Current Named Example: Bain}

PyTex currently exposes one named helper, the Bain correspondence constructor for cubic parent and
child phases. In the present implementation this helper encodes

\begin{align}
(001)_p &\parallel (001)_c \\
[110]_p &\parallel [100]_c
\end{align}

and then delegates to the same plane-direction basis construction described above.

\subsection{Current Named Example: Nishiyama--Wassermann}

PyTex also exposes a Nishiyama--Wassermann helper for cubic parent and child phases. In the present
implementation this helper encodes

\begin{align}
(111)_p &\parallel (011)_c \\
[1\bar{1}0]_p &\parallel [100]_c
\end{align}

and again delegates to the same plane-direction basis construction described above.
